% Generated draft: requires mathematical and source review.
\documentclass[11pt]{article}
\usepackage[margin=25mm]{geometry}
\usepackage{fontspec}
\setmainfont{Latin Modern Roman}
\usepackage{amsmath,amssymb,mathtools,unicode-math}
\setmathfont{Latin Modern Math}
\usepackage{enumitem,needspace,xurl,hyperref}
\setlength{\parindent}{0pt}
\setlength{\parskip}{5pt}
\setlength{\emergencystretch}{3em}
\newcommand{\N}{\mathbb{N}}
\newcommand{\Z}{\mathbb{Z}}
\newcommand{\Q}{\mathbb{Q}}
\newcommand{\R}{\mathbb{R}}
\newcommand{\C}{\mathbb{C}}
\newcommand{\F}{\mathbb{F}}
\newcommand{\A}{\mathbb{A}}
\newcommand{\PP}{\mathbb{P}}
\newcommand{\E}{\mathbb{E}}
\newcommand{\eps}{\varepsilon}
\newcommand{\dd}{\,\mathrm{d}}
\begin{document}
\section*{UnsolvedMath ID: 5500026}
\textbf{Source problem code:} catalogue entry
\subsection{Definitions and mathematical statement}
In the smooth closed-surface setting described by the supplied research summary, let \(S\subseteq\R^3\) be a compact embedded smooth surface without boundary. Let \(\operatorname{Med}(S)\) be its medial axis, using the closure of points having more than one nearest point on \(S\), and set
\[
\operatorname{lfs}_S(x)=\operatorname{dist}(x,\operatorname{Med}(S))\quad(x\in S).
\]
A finite noiseless sample \(P\subseteq S\) is an \(\varepsilon\)-sample in the customary local-feature-size sense when
\[
\forall x\in S,\qquad \operatorname{dist}(x,P)\le
\varepsilon\operatorname{lfs}_S(x).
\]
Thus density is relative to local geometry, not just an absolute spacing bound.

The reconstruction task is to give an algorithm using the sample as input that, under an appropriate sufficiently small \(\varepsilon\) guarantee, produces a surface \(\widehat S(P)\) and guarantees a homeomorphism
\[
h:S\longrightarrow\widehat S(P).
\]
In the usual finite-output formulation \(\widehat S(P)\) is a triangulated polyhedral surface. The theorem must specify the admissible class of \(S\), a sufficient sampling threshold, and the algorithm's input and output model. It need not recover exact geometry or certify that the input satisfies an unknown source surface's sampling condition.

The supplied research summary reports provable reconstruction algorithms called Crust, Cocone, and Power Crust for smooth closed surfaces under sampling assumptions. No numerical threshold or exact theorem hypotheses are supplied, so the algorithms are not asserted here to share identical assumptions. The broad original wording merely says surface; it does not itself state smoothness, closedness, ambient dimension, treatment of boundaries or singularities, or noise assumptions. The smooth closed embedded \(\R^3\) framework above makes the reported solved scope explicit rather than silently extending it to arbitrary surfaces.
\subsection{Short English statement}
Reconstruct a surface homeomorphic to an unknown surface from a sample dense relative to its local feature size. The supplied literature assessment reports solutions for smooth closed surfaces, while the broader source wording needs its exact assumptions fixed.
\subsection{Sources}
\begin{itemize}
\item Local UnsolvedMath registry, permanent ID \(5500026\), exact source code \texttt{catalogue entry}, ``Surface Reconstruction.''
\item \emph{The Open Problems Project}, Computational Geometry, Problem 26; \url{https://topp.openproblem.net/p26}. Source-TeX extraction; recorded access 2026-07-29. The recorded maintained-entry status says open.
\item Supplied research summary names Crust, Cocone, and Power Crust as provable reconstruction algorithms, but gives no primary publications, authors, precise sampling constants, or theorem statements.
\item Source provenance repository: \url{https://github.com/edemaine/topp}; the recorded license expressly scopes software, leaving problem-text reuse for release review.
\end{itemize}
\end{document}
